\section*{Bab \ref{ch:b3:bioinformatics} --- Bioinformatika dan Analisis Urutan}

\begin{solution}{exo:b3:bioinformatics:1}
Global: kedua urutannya dijajarkan dari ujung ke ujung, tiap residu berada di
dalam sebuah lajur --- bagi dua protein yang diyakini homolog sepanjang
seluruhnya, misalnya \omterm{def:b3:genomics:comparative}{ortolog} sebuah enzim rumah tangga. Lokal: pasangan
subuntai berskor terbaik, sisanya diabaikan --- bagi pencarian domain bersama
(sebuah domain SH2 pada dua protein pengisyarat yang selebihnya tak
berkerabat), atau sebuah gen di dalam urutan genom yang panjang.
\end{solution}

\begin{solution}{exo:b3:bioinformatics:2}
Batasnya $0,-1,-2,-3$ pada kedua arah. Baris A: $1, 0, -1$. Baris G: $0, 0, -1$.
Baris C: $-1, -1, 1$. Optimumnya $F(3,3) = 1$: AGC di atas AAC tanpa celah
(cocok, tak cocok, cocok: $1 - 1 + 1 = 1$).
\end{solution}

\begin{solution}{exo:b3:bioinformatics:3}
$s(a,a) = \lambda^{-1}\log\bigl(q_{aa}/p_{a}^{2}\bigr)$. Triptofan itu
langka ($p_{W} \approx 0.013$), sehingga peluang dua triptofan terjajar
secara acak, $p_{W}^{2}$, sangat kecil, dan pasangan triptofan yang lestari
menjadi tanda homologi yang jauh lebih kuat daripada pasangan leusina yang
lestari ($p_{L}
\approx 0.1$); nisbah log-odsnya pun sebanding lebih besar.
\end{solution}

\begin{solution}{exo:b3:bioinformatics:4}
Nilai-E adalah cacah \omterm{def:b3:bioinformatics:alignment}{penjajaran} berskor sedikitnya setinggi itu
yang akan diharapkan muncul secara kebetulan pada penelusuran kueri ini
terhadap pangkalan data seukuran ini. $E = 3$ berarti tiga skor semacam itu
diharapkan muncul secara kebetulan: kecocokannya bukan bukti homologi (boleh
jadi ia memang homolog, tetapi penelusurannya tak dapat memutuskan).
\end{solution}

\begin{solution}{exo:b3:bioinformatics:5}
$mn = 400\times 2\times 10^{11} = 8\times 10^{13} = 2^{46.2}$. $E(45) =
2^{1.2} \approx 2.3$; $E(55) = 2^{-8.8} \approx 2\times 10^{-3}$;
$E(65) = 2^{-18.8} \approx 2\times 10^{-6}$. $E = 10^{-3}$ menuntut $S' =
46.2 + 10.0 = 56$ bit. Kueri berisi \num{40} residu memiliki $mn$ sepuluh kali
lebih kecil, $2^{42.9}$: $53$ bit memadai --- tetapi kueri pendek jarang
mencapai bahkan angka itu.
\end{solution}

\begin{solution}{exo:b3:bioinformatics:6}
\omterm{def:b3:bioinformatics:motif}{Kandungan informasinya}: $2$, $1$, $0$, dan $2 - H$ dengan $H = -(0.7\log_{2}
0.7 + 3\times 0.1\log_{2} 0.1) = 0.36 + 1.00 = 1.36$, jadi $0.64$. Totalnya
$R = 3.64$ bit. Kecocokan kebetulannya: $9.2\times 10^{6}$ kedudukan pada dua
untai $\times 2^{-3.64} \approx 7\times 10^{5}$ --- \omterm{def:b3:bioinformatics:motif}{motifnya} nyaris
tak berguna sendirian.
\end{solution}

\begin{solution}{exo:b3:bioinformatics:7}
Sebuah sisipan berisi beberapa residu merupakan satu peristiwa mutasi,
sehingga ongkosnya tak semestinya tumbuh lurus dengan panjangnya: ongkos
pembukaan ditambah ongkos perpanjangan yang kecil memodelkan hal itu. Denda
pembukaan yang terlalu tinggi memaksa ketakcocokan di tempat sebuah celah
semestinya berada lalu salah menjajarkan segala sesuatu sesudah sebuah sisipan
yang sungguhan; denda yang terlalu rendah menaburkan celah di mana-mana,
memasangkan residu secara kebetulan dan menggelembungkan keidentikannya.
\end{solution}

\begin{solution}{exo:b3:bioinformatics:8}
Belum tentu: \omterm{def:b3:bioinformatics:alignment}{penjajarannya} meliputi segmen berisi \num{130} residu,
yang merupakan ciri khas domain bersama, bukan ciri \omterm{def:b3:genomics:comparative}{ortolog}
yang terjajar sepanjang seluruhnya. Ujinya: telusurkanlah protein lalatnya
kembali pada proteom manusia (apakah kuerinya menjadi kecocokan terbaiknya,
sepanjang seluruhnya?), kenalilah domainnya dengan HMM \omterm{def:b3:bioinformatics:msa}{profil}, lalu bangunlah
sebuah pohon gen famili itu pada beberapa spesies.
\end{solution}

\begin{solution}{exo:b3:bioinformatics:9}
Sebuah \omterm{def:b3:bioinformatics:msa}{profil} mencatat, lajur demi lajur, apa yang ditoleransi familinya:
sebuah kedudukan yang selalu hidrofob tetapi tak pernah residu yang sama,
sebuah residu katalitik yang tetap, sebuah kedudukan yang selalu berupa celah
pada separuh familinya. \omterm{def:b3:bioinformatics:alignment}{Penjajaran} berpasangan memberi skor tiap residu
terhadap satu residu lain saja dan tak dapat mengetahui bahwa sebuah valina
pada kedudukan 40 ``sama baiknya dengan'' isoleusina yang ada di situ pada
kuerinya. \omterm{def:b3:bioinformatics:msa}{Profilnya} juga membobot lajur yang lestari, sehingga kemiripan lemah
yang terpusat di tempat familinya lestari menjadi bermakna.
\end{solution}

\begin{solution}{exo:b3:bioinformatics:10}
Dengan \omterm{prop:b3:bioinformatics:logodds}{harapan skor} $\mu > 0$ tiap lajur, skor kumulatif
sepanjang diagonal dua urutan acak merupakan jalan acak berhanyutan
positif: sesudah $n$ lajur skornya kira-kira $\mu n$, sehingga \omterm{def:b3:bioinformatics:alignment}{penjajaran
lokal} terbaiknya pada dasarnya seluruhnya dan skornya tumbuh sebagai
$\mu n$, bukan sebagai $\log n$. Teori Karlin--Altschul, yang
menuntut hanyutan negatif agar skor tinggi menjadi penyimpangan yang langka,
tak berlaku dan tak ada $\lambda$ yang ada. Bagi DNA pada \qty{60}{\%} GC,
peluang sebuah kecocokan adalah $2(0.3^{2}) + 2(0.2^{2}) = 0.26$, sehingga
\omterm{prop:b3:bioinformatics:logodds}{harapan skornya} $0.26 - 0.74 = -0.48$: masih negatif, dan statistikanya
berlaku; tetapi skema semacam kecocokan $+1$, ketakcocokan $-0.3$ akan
berharapan $+0.04$ dan akan melaporkan seluruh genomnya sebagai satu \omterm{def:b3:bioinformatics:alignment}{penjajaran}.
\end{solution}

\begin{solution}{exo:b3:bioinformatics:11}
\omterm{def:b3:bioinformatics:blast}{Benih} dan perangkaian: carilah kecocokan $k$-mer yang persis atau hampir
persis antara kedua urutannya dengan tabel hash, simpanlah yang berbaris pada
diagonal yang taat asas, rangkaikanlah, lalu jalankanlah \omterm{thm:b3:bioinformatics:nw}{pemrograman dinamis}
hanya pada celah di antara \omterm{def:b3:bioinformatics:blast}{benih} yang terangkai; cara itu mengorbankan
\omterm{def:b3:bioinformatics:alignment}{penjajaran} di daerah yang tanpa \omterm{def:b3:bioinformatics:blast}{benih} (bentangan yang sangat menyimpang).
Pemitaan: bila kedua urutannya diketahui nyaris segaris, hitunglah hanya sel
di dalam pita selebar $w$ di sekitar diagonalnya, dengan biaya $wn$ alih-alih
$mn$; cara itu mengorbankan \omterm{def:b3:bioinformatics:alignment}{penjajaran} mana pun yang sisipannya lebih besar
daripada pitanya.
\end{solution}

\begin{solution}{exo:b3:bioinformatics:12}
Urutan penyandi berperiode tiga: ketiga kedudukan kodonnya memiliki
komposisi basa yang berbeda (kedudukan ketiganya paling bias), dan pemakaian
kodonnya tak merata pada tiap spesies. Sebuah model dengan tiga keadaan
penyandi yang berurutan, yang masing-masing memancarkan basa dengan komposisi
kedudukan kodon itu, memberi DNA penyandi peluang yang lebih tinggi daripada
yang diberikan keadaan tak penyandi, sepanjang jendela beberapa lusin kodon,
bahkan tanpa kodon henti. Di dalam genom manusia eksonnya pendek
(\qty{150}{bp}), terpisah oleh intron sepanjang kilobasa, sehingga isyarat
penyandinya singkat dan terputus; modelnya harus pula mengenali tapak sambung,
yang isyaratnya lemah, dan urutan tak penyandi yang begitu banyak menghasilkan
banyak segmen penyandi palsu.
\end{solution}

\begin{solution}{pb:b3:bioinformatics:1}
\textbf{1.} Baris K, Q, T; lajur K, A, Q, T; batasnya $0,-1,-2,-3,-4$
dan $0,-1,-2,-3$. Baris K: $1, 0, -1, -2$; baris Q: $0, 0, 1, 0$; baris T:
$-1, -1, 0, 2$. Optimumnya $2$: \texttt{K-QT} di atas \texttt{KAQT}.
\textbf{2.} Skor lokal terbaiknya $3$: \texttt{CAT} terhadap \texttt{CAT}
(residu 4--6 GATCAT dengan 2--4 ACAT); \texttt{ATCAT} terhadap
\texttt{A-CAT} pun berskor $4 - 1 = 3$.
\textbf{3.} $300\times 450 = 1.35\times 10^{5}$ pembaruan; terhadap
pangkalan datanya $300\times 1.2\times 10^{11} = 3.6\times 10^{13}$.
\textbf{4.} $3.6\times 10^{4}$ s, sepuluh jam tiap kueri; \omterm{def:b3:bioinformatics:blast}{benih} \omterm{def:b3:bioinformatics:blast}{BLAST}
melompati hampir seluruh tabelnya dan menjawab dalam hitungan detik.
\textbf{5.} $q_{ab} = p_{a}p_{b}e^{\lambda s}$. Alanina: $e^{1.39} = 4.0$,
$q_{AA} = 0.074^{2}\times 4.0 = 0.022$, nisbahnya $4$. Triptofan:
$e^{3.82} = 45$, $q_{WW} = 0.013^{2}\times 45 = 0.0077$, nisbahnya $45$.
Sepasang triptofan yang terjajar $45$ kali lebih sering pada homolog daripada
secara kebetulan, sepasang alanina hanya empat kali; namun pasangan alanina
lebih banyak secara mutlak karena alanina memang lazim.
\textbf{6.} \qty{24}{\%} terletak di zona remang, tempat \omterm{def:b3:bioinformatics:alignment}{penjajaran}
acak mencapai \qtyrange{15}{20}{\%}; nilai-E \omterm{def:b3:bioinformatics:alignment}{penjajarannya},
\omterm{def:b3:bioinformatics:motif}{motif} lestari pada kedudukan yang tepat, kecocokan HMM \omterm{def:b3:bioinformatics:msa}{profil} dengan famili
yang dikenal, atau lipatan bersama akan memutuskannya.
\textbf{7.} $mn = 300\times 1.2\times 10^{11} = 3.6\times 10^{13}$;
$\log_{2}(mn) = 45.0$.
\textbf{8.} $E(92) = 2^{45 - 92} = 2^{-47} \approx 7\times 10^{-15}$;
$E(38) = 2^{7} = 128$.
\textbf{9.} $E = 10^{-3}$ pada $S' = 45 + 10 = 55$ bit; $E = 1$ pada $45$
bit.
\textbf{10.} Sepuluh kali $mn$: $E \approx 7\times 10^{-14}$, masih
sangat meyakinkan.
\textbf{11.} Dengan $E = 128$, kecocokan kesepuluhnya adalah apa yang
dihasilkan kebetulan; keidentikan \qty{40}{\%} sepanjang $60$ residu bukan
bukti. Penelusuran \omterm{def:b3:bioinformatics:msa}{profil} atas residu 200--260 pada pangkalan data domain
dapat memperlihatkan apakah segmen itu sebuah domain yang dikenal, dengan
statistika yang tak dimiliki perbandingan berpasangan.
\textbf{12.} Kecocokan terbaik timbal balik sepanjang seluruhnya selaras
dengan ortologi satu lawan satu; keduanya tidak membuktikannya --- sebuah
penggandaan pada satu garis keturunan sesudah perpisahannya memberi dua
koortolog, dan hilangnya \omterm{def:b3:genomics:comparative}{ortolog} yang sungguhan dapat meninggalkan sebuah
\omterm{def:b3:genomics:comparative}{paralog} sebagai kecocokan terbaiknya. Sebuah pohon gen dengan beberapa
spesies adalah ujinya.
\textbf{13.} $R = 2 + 2 + 1.6 + 2 + 0.8 + 1.2 + 0.4 + 0.3 = 10.3$ bit.
\textbf{14.} $3.2\times 10^{9}\times 2^{-10.3} \approx 2.5\times 10^{6}$
kecocokan kebetulan.
\textbf{15.} $\log_{2}(3.2\times 10^{9}/200) = \log_{2}(1.6\times 10^{7})
\approx 24$ bit.
\textbf{16.} Kemunculan bersama pada jarak tetap menambahkan bitnya: $10.3 + 8 =
18.3$, yang memasok $8$ dari $13.7$ bit yang hilang; kira-kira $5.7$ bit
(faktor $50$ pada kecocokan kebetulan) harus datang dari tempat lain ---
keterjangkauan \omterm{def:b3:chromatin-epigenetics:nucleosome}{kromatin}, mitra yang lain.
\textbf{17.} $H = -(0.5\log_{2}0.5 + 0.5\log_{2}0.5) = 1$ bit; $R = 2 -
1 = 1$ bit.
\textbf{18.} Sebuah faktor bakteri harus menemukan tapaknya yang sedikit di
dalam genom \qty{4.6}{Mb} tanpa bantuan \omterm{def:b3:chromatin-epigenetics:nucleosome}{kromatin}: ia menuntut sekitar $19$
bit, dan tapaknya panjang serta lestari. Genom eukariot seribu kali lebih
besar, sehingga menuntut sepuluh bit lagi, namun faktornya bertapak pendek;
kekhasannya dicapai lewat gabungan dan lewat pembatasan pada \omterm{def:b3:chromatin-epigenetics:nucleosome}{kromatin} yang
terjangkau, yang juga membuat pengaturannya lebih mudah berevolusi, karena
tapak pendek mudah diperoleh atau hilang.
\textbf{19.} $3/64 = 0.047$; harapan cacah kodon sebelum sebuah kodon henti
adalah $64/3 \approx 21$.
\textbf{20.} $p_{A} = p_{T} = 0.31$, $p_{G} = p_{C} = 0.19$: $P(\text{TAA})
= 0.31^{3} = 0.030$, $P(\text{TAG}) = P(\text{TGA}) = 0.31^{2}\times 0.19
= 0.018$; totalnya $0.066$, harapan panjang kerangkanya $15$ kodon. DNA yang
kaya AT penuh kodon henti, sehingga kerangka terbuka acak lebih pendek dan
kerangka yang panjang lebih menonjol.
\textbf{21.} $(61/64)^{100} = e^{-4.80} = 0.008$; $(61/64)^{300} =
e^{-14.4} = 5.6\times 10^{-7}$.
\textbf{22.} Tiap kerangka terbuka maksimal berakhir pada sebuah kodon henti,
dan $3.2\times
10^{9}$ awalan kodon memuat $3.2\times 10^{9}\times 3/64 = 1.5\times
10^{8}$ kodon henti: kira-kira $1.5\times 10^{8}\times 0.008 = 1.2\times 10^{6}$
kerangka acak sepanjang sedikitnya $100$ kodon, dan $1.5\times 10^{8}\times
5.6\times 10^{-7} \approx 80$ yang sedikitnya $300$.
\textbf{23.} Sebuah bakteri \qty{4.6}{Mb} memiliki sekitar $4\times 10^{5}$
kodon henti dan karenanya sekitar \num{3500} kerangka kebetulan sepanjang $100$
kodon tetapi nyaris tak satu pun sepanjang $300$; gennya rata-rata $300$ kodon
dan \qty{88}{\%} DNA-nya penyandi, sehingga kerangka terbuka yang panjang
hampir selalu sebuah gen. Pada cacing itu, \qty{1.5}{\%} DNA-nya menyandi,
eksonnya rata-rata $50$ kodon --- lebih pendek daripada ambang kebetulannya ---
dan sejuta kerangka acak sepanjang $100$ kodon menenggelamkannya. Pencari gen
eukariot memakai isyarat tapak sambung, bias kodon di dalam model Markov
tersembunyi, homologi dengan protein yang dikenal dan, di atas segalanya,
transkrip yang telah diurutkan.
\textbf{24.} Sebuah \omterm{def:b3:genomics:ngs}{bacaan} dari RNA duta yang tersambung terjajar pada
genomnya menjadi dua keping yang dipisahkan sebuah intron: belahannya menandai
kedua tapak sambungnya sampai ke basanya; cakupan \omterm{def:b3:genomics:ngs}{bacaannya} membatasi
eksonnya dan \omterm{def:b3:genomics:ngs}{bacaan} berpasangan menghubungkan ekson berurutan menjadi satu
transkrip, sehingga terurai tapak sambung calon mana yang terpakai.
\textbf{25.} $E \approx 7\times 10^{-15}$ bagi kecocokan terbaiknya; sekitar
$24$ bit untuk menentukan $200$ tapak di dalam genomnya; sekitar $80$ kerangka
baca kebetulan sepanjang $300$ kodon di dalam seluruh genomnya.
\end{solution}
